\documentclass[12pt,a4paper]{article}
\usepackage[UTF8]{ctex}
\usepackage{amsmath,amssymb,amsthm}
\usepackage{geometry}

\geometry{margin=2.5cm}

\title{三个动力学方程的理解}
\author{第8章补充说明}
\date{}

\begin{document}

\maketitle

\section{问题提出}

本文档解释三个相关的动力学方程：
\begin{enumerate}
\item $\tau = M(\theta)\ddot{\theta} + h(\theta, \dot{\theta})$ \quad (8.82)
\item $F = \Lambda(\theta)\dot{V} + \eta(\theta, V)$ \quad (8.89)
\item $M(\theta)\ddot{\theta} + h(\theta, \dot{\theta}) - J^T(\theta)F_{\text{ext}}$ \quad (1)
\end{enumerate}

\section{公式1：关节空间动力学（无外部力）}

\subsection{公式}

\begin{equation}
\tau = M(\theta)\ddot{\theta} + h(\theta, \dot{\theta}) \tag{8.82}
\end{equation}

\subsection{各项含义}

\begin{itemize}
\item $\tau \in \mathbb{R}^n$：\textbf{关节力矩向量}（由执行器提供）
\item $M(\theta) \in \mathbb{R}^{n \times n}$：\textbf{关节空间质量矩阵}（对称正定）
\item $\ddot{\theta} \in \mathbb{R}^n$：\textbf{关节加速度向量}
\item $h(\theta, \dot{\theta}) \in \mathbb{R}^n$：\textbf{科里奥利、向心力和重力项}
\end{itemize}

\subsection{物理意义}

\textbf{这个方程描述}：
\begin{itemize}
\item 要产生关节加速度$\ddot{\theta}$，需要施加关节力矩$\tau$
\item $\tau$包括两部分：
\begin{enumerate}
\item 惯性力：$M(\theta)\ddot{\theta}$（产生加速度所需的力）
\item 科里奥利、向心力和重力：$h(\theta, \dot{\theta})$（由于速度和位置产生的力）
\end{enumerate}
\end{itemize}

\subsection{假设条件}

\begin{itemize}
\item \textbf{无外部力}：假设末端执行器上没有外部力作用
\item \textbf{理想情况}：忽略摩擦、间隙等非理想因素
\end{itemize}

\subsection{例子：单关节机器人}

\textbf{单关节动力学}：
\begin{equation}
\tau = M\ddot{\theta} + mgr\cos\theta \tag{11.19}
\end{equation}

其中：
\begin{itemize}
\item $M$：连杆关于旋转轴的标量惯性
\item $m$：连杆质量
\item $r$：从轴到质心的距离
\item $g$：重力加速度
\end{itemize}

\textbf{分解}：
\begin{itemize}
\item $M\ddot{\theta}$：惯性力（产生加速度所需的力矩）
\item $mgr\cos\theta$：重力项（依赖于角度$\theta$）
\end{itemize}

\section{公式3：关节空间动力学（含外部力）}

\subsection{公式}

\begin{equation}
\tau = M(\theta)\ddot{\theta} + h(\theta, \dot{\theta}) - J^T(\theta)F_{\text{ext}} \tag{1}
\end{equation}

\subsection{与公式1的对比}

\textbf{公式1}（无外部力）：
\begin{equation}
\tau = M(\theta)\ddot{\theta} + h(\theta, \dot{\theta})
\end{equation}

\textbf{公式3}（含外部力）：
\begin{equation}
\tau = M(\theta)\ddot{\theta} + h(\theta, \dot{\theta}) - J^T(\theta)F_{\text{ext}}
\end{equation}

\textbf{区别}：
\begin{itemize}
\item 公式3多了一项：$-J^T(\theta)F_{\text{ext}}$
\item 这一项表示外部力$F_{\text{ext}}$对关节力矩的影响
\end{itemize}

\subsection{各项含义}

\begin{itemize}
\item $\tau$：关节力矩（由执行器提供）
\item $M(\theta)\ddot{\theta}$：惯性力
\item $h(\theta, \dot{\theta})$：科里奥利、向心力和重力
\item $F_{\text{ext}} \in \mathbb{R}^6$：\textbf{作用在末端执行器上的外部wrench}（6×1向量，包含力和力矩）
\item $J^T(\theta)$：\textbf{雅可比矩阵的转置}（将任务空间wrench转换为关节空间力矩）
\item $-J^T(\theta)F_{\text{ext}}$：\textbf{外部力在关节空间中的等效力矩}
\end{itemize}

\subsection{物理意义}

\textbf{这个方程描述}：
\begin{itemize}
\item 要产生关节加速度$\ddot{\theta}$，需要施加关节力矩$\tau$
\item 但外部力$F_{\text{ext}}$会"帮助"或"阻碍"运动
\item 如果$F_{\text{ext}}$与运动方向相同，$J^T(\theta)F_{\text{ext}}$为正，需要的$\tau$减小
\item 如果$F_{\text{ext}}$与运动方向相反，$J^T(\theta)F_{\text{ext}}$为负，需要的$\tau$增大
\end{itemize}

\textbf{重新整理}：
\begin{equation}
\tau + J^T(\theta)F_{\text{ext}} = M(\theta)\ddot{\theta} + h(\theta, \dot{\theta})
\end{equation}

\textbf{物理意义}：
\begin{itemize}
\item 总力矩（执行器力矩 + 外部力等效力矩）= 惯性力 + 其他力
\end{itemize}

\subsection{虚功原理}

\textbf{为什么是$J^T(\theta)$？}

\textbf{虚功原理}告诉我们，关节空间和任务空间之间的功率是守恒的：
\begin{equation}
\tau^T\dot{\theta} = F^T V
\end{equation}

其中$V = J(\theta)\dot{\theta}$是末端执行器的twist。

\textbf{推导}：
\begin{align}
\tau^T\dot{\theta} &= F^T V \\
&= F^T J(\theta)\dot{\theta} \\
&= (J^T(\theta)F)^T\dot{\theta}
\end{align}

因此：
\begin{equation}
\tau = J^T(\theta)F
\end{equation}

\textbf{物理意义}：
\begin{itemize}
\item 任务空间wrench $F$在关节空间中的等效力矩是$J^T(\theta)F$
\item 如果$F$是外部力，那么$-J^T(\theta)F$表示外部力对关节力矩的"贡献"
\end{itemize}

\section{公式2：任务空间动力学}

\subsection{公式}

\begin{equation}
F = \Lambda(\theta)\dot{V} + \eta(\theta, V) \tag{8.89}
\end{equation}

\subsection{各项含义}

\begin{itemize}
\item $F \in \mathbb{R}^6$：\textbf{任务空间wrench}（6×1向量，包含力和力矩）
\item $\Lambda(\theta) \in \mathbb{R}^{6 \times 6}$：\textbf{任务空间质量矩阵}（对称正定）
\item $\dot{V} \in \mathbb{R}^6$：\textbf{末端执行器加速度}（twist的导数）
\item $\eta(\theta, V) \in \mathbb{R}^6$：\textbf{任务空间的科里奥利、向心力和重力项}
\end{itemize}

\subsection{物理意义}

\textbf{这个方程描述}：
\begin{itemize}
\item 要产生末端执行器加速度$\dot{V}$，需要施加wrench $F$
\item $F$包括两部分：
\begin{enumerate}
\item 惯性力：$\Lambda(\theta)\dot{V}$（产生加速度所需的力）
\item 科里奥利、向心力和重力：$\eta(\theta, V)$（由于速度和位置产生的力）
\end{enumerate}
\end{itemize}

\subsection{与公式1和公式3的关系}

\textbf{公式2是从公式1或公式3转换而来的}：

\textbf{从公式1转换}（假设$\tau = J^T(\theta)F$，无外部力）：
\begin{equation}
J^T(\theta)F = M(\theta)\ddot{\theta} + h(\theta, \dot{\theta})
\end{equation}

左乘$J^{-T}(\theta)$并转换，得到：
\begin{equation}
F = \Lambda(\theta)\dot{V} + \eta(\theta, V)
\end{equation}

\textbf{从公式3转换}（含外部力）：
\begin{equation}
J^T(\theta)\tau + F_{\text{ext}} = M(\theta)\ddot{\theta} + h(\theta, \dot{\theta})
\end{equation}

转换后：
\begin{equation}
F_{\text{total}} = \Lambda(\theta)\dot{V} + \eta(\theta, V)
\end{equation}

其中$F_{\text{total}} = J^T(\theta)\tau + F_{\text{ext}}$是总wrench。

\subsection{转换关系}

\textbf{质量矩阵}：
\begin{equation}
\Lambda(\theta) = J^{-T}(\theta)M(\theta)J^{-1}(\theta) \tag{8.90}
\end{equation}

\textbf{科里奥利等项}：
\begin{equation}
\eta(\theta, V) = J^{-T}(\theta)h(\theta, J^{-1}(\theta)V) - \Lambda(\theta)\dot{J}(\theta)J^{-1}(\theta)V \tag{8.91}
\end{equation}

\textbf{速度关系}：
\begin{align}
V &= J(\theta)\dot{\theta} \tag{8.83} \\
\dot{V} &= \dot{J}(\theta)\dot{\theta} + J(\theta)\ddot{\theta} \tag{8.84}
\end{align}

\section{三个公式的对比总结}

\subsection{公式对比表}

\begin{center}
\begin{tabular}{|l|l|l|l|}
\hline
\textbf{公式} & \textbf{坐标系} & \textbf{变量} & \textbf{是否含外部力} \\
\hline
(8.82) & 关节空间 & $\tau, \theta, \dot{\theta}, \ddot{\theta}$ & 否 \\
\hline
(1) & 关节空间 & $\tau, \theta, \dot{\theta}, \ddot{\theta}, F_{\text{ext}}$ & 是 \\
\hline
(8.89) & 任务空间 & $F, V, \dot{V}$ & 取决于$F$的含义 \\
\hline
\end{tabular}
\end{center}

\subsection{公式1：$\tau = M(\theta)\ddot{\theta} + h(\theta, \dot{\theta})$}

\textbf{特点}：
\begin{itemize}
\item 关节空间表示
\item 无外部力
\item 适用于自由空间运动
\end{itemize}

\textbf{应用}：
\begin{itemize}
\item 自由空间轨迹跟踪
\item 无接触任务
\end{itemize}

\subsection{公式3：$\tau = M(\theta)\ddot{\theta} + h(\theta, \dot{\theta}) - J^T(\theta)F_{\text{ext}}$}

\textbf{特点}：
\begin{itemize}
\item 关节空间表示
\item 含外部力
\item 适用于接触任务
\end{itemize}

\textbf{应用}：
\begin{itemize}
\item 力控制
\item 接触任务（装配、打磨等）
\item 阻抗控制
\end{itemize}

\subsection{公式2：$F = \Lambda(\theta)\dot{V} + \eta(\theta, V)$}

\textbf{特点}：
\begin{itemize}
\item 任务空间表示
\item 直接描述末端执行器动力学
\item 适用于任务空间控制
\end{itemize}

\textbf{应用}：
\begin{itemize}
\item 任务空间轨迹跟踪
\item 任务空间力控制
\item 任务空间阻抗控制
\end{itemize}

\section{实际应用场景}

\subsection{场景1：自由空间运动}

\textbf{使用公式1}：
\begin{equation}
\tau = M(\theta)\ddot{\theta} + h(\theta, \dot{\theta})
\end{equation}

\textbf{控制律}：
\begin{equation}
\tau = M(\theta)\ddot{\theta}_d + h(\theta, \dot{\theta}) + K_p(\theta_d - \theta) + K_d(\dot{\theta}_d - \dot{\theta})
\end{equation}

\textbf{特点}：
\begin{itemize}
\item 无外部力
\item 直接控制关节
\end{itemize}

\subsection{场景2：接触任务}

\textbf{使用公式3}：
\begin{equation}
\tau = M(\theta)\ddot{\theta} + h(\theta, \dot{\theta}) - J^T(\theta)F_{\text{ext}}
\end{equation}

\textbf{控制律}（力控制）：
\begin{equation}
\tau = M(\theta)\ddot{\theta}_d + h(\theta, \dot{\theta}) - J^T(\theta)(F_d + K_f(F_d - F_{\text{ext}}))
\end{equation}

\textbf{特点}：
\begin{itemize}
\item 考虑外部力
\item 可以控制接触力
\end{itemize}

\subsection{场景3：任务空间控制}

\textbf{使用公式2}：
\begin{equation}
F = \Lambda(\theta)\dot{V} + \eta(\theta, V)
\end{equation}

\textbf{控制律}：
\begin{equation}
F = \Lambda(\theta)\dot{V}_d + \eta(\theta, V) + K_p X_e + K_d V_e
\end{equation}

\textbf{转换到关节空间}：
\begin{equation}
\tau = J^T(\theta)F
\end{equation}

\textbf{特点}：
\begin{itemize}
\item 任务空间表示
\item 直接控制末端执行器
\end{itemize}

\section{数学推导：从公式1到公式2}

\subsection{步骤1：从公式1开始}

\begin{equation}
\tau = M(\theta)\ddot{\theta} + h(\theta, \dot{\theta}) \tag{8.82}
\end{equation}

\subsection{步骤2：假设$\tau = J^T(\theta)F$}

这表示关节力矩$\tau$是由任务空间wrench $F$产生的。

\begin{equation}
J^T(\theta)F = M(\theta)\ddot{\theta} + h(\theta, \dot{\theta})
\end{equation}

\subsection{步骤3：左乘$J^{-T}(\theta)$}

\begin{equation}
F = J^{-T}(\theta)M(\theta)\ddot{\theta} + J^{-T}(\theta)h(\theta, \dot{\theta})
\end{equation}

\subsection{步骤4：代入速度关系}

从$V = J(\theta)\dot{\theta}$，得到：
\begin{equation}
\dot{\theta} = J^{-1}(\theta)V
\end{equation}

从$\dot{V} = \dot{J}(\theta)\dot{\theta} + J(\theta)\ddot{\theta}$，得到：
\begin{equation}
\ddot{\theta} = J^{-1}(\theta)(\dot{V} - \dot{J}(\theta)J^{-1}(\theta)V)
\end{equation}

\subsection{步骤5：代入并整理}

\begin{align}
F &= J^{-T}(\theta)M(\theta)J^{-1}(\theta)(\dot{V} - \dot{J}(\theta)J^{-1}(\theta)V) + J^{-T}(\theta)h(\theta, J^{-1}(\theta)V) \\
&= J^{-T}(\theta)M(\theta)J^{-1}(\theta)\dot{V} \\
&\quad - J^{-T}(\theta)M(\theta)J^{-1}(\theta)\dot{J}(\theta)J^{-1}(\theta)V \\
&\quad + J^{-T}(\theta)h(\theta, J^{-1}(\theta)V)
\end{align}

\subsection{步骤6：定义$\Lambda(\theta)$和$\eta(\theta, V)$}

\begin{align}
\Lambda(\theta) &= J^{-T}(\theta)M(\theta)J^{-1}(\theta) \tag{8.90} \\
\eta(\theta, V) &= J^{-T}(\theta)h(\theta, J^{-1}(\theta)V) - \Lambda(\theta)\dot{J}(\theta)J^{-1}(\theta)V \tag{8.91}
\end{align}

\subsection{步骤7：得到最终方程}

\begin{equation}
F = \Lambda(\theta)\dot{V} + \eta(\theta, V) \tag{8.89}
\end{equation}

\section{数学推导：从公式3到公式2}

\subsection{步骤1：从公式3开始}

\begin{equation}
\tau = M(\theta)\ddot{\theta} + h(\theta, \dot{\theta}) - J^T(\theta)F_{\text{ext}} \tag{1}
\end{equation}

\subsection{步骤2：重新整理}

\begin{equation}
\tau + J^T(\theta)F_{\text{ext}} = M(\theta)\ddot{\theta} + h(\theta, \dot{\theta})
\end{equation}

\subsection{步骤3：假设$\tau = J^T(\theta)F_{\text{actuator}}$}

这表示关节力矩$\tau$是由执行器产生的任务空间wrench $F_{\text{actuator}}$产生的。

\begin{equation}
J^T(\theta)(F_{\text{actuator}} + F_{\text{ext}}) = M(\theta)\ddot{\theta} + h(\theta, \dot{\theta})
\end{equation}

\subsection{步骤4：定义总wrench}

\begin{equation}
F_{\text{total}} = F_{\text{actuator}} + F_{\text{ext}}
\end{equation}

\subsection{步骤5：转换到任务空间}

按照从公式1到公式2的步骤，得到：
\begin{equation}
F_{\text{total}} = \Lambda(\theta)\dot{V} + \eta(\theta, V)
\end{equation}

\textbf{如果$F_{\text{ext}} = 0$}（无外部力），则：
\begin{equation}
F_{\text{actuator}} = \Lambda(\theta)\dot{V} + \eta(\theta, V)
\end{equation}

这就是公式2（其中$F = F_{\text{actuator}}$）。

\section{总结}

\subsection{三个公式的关系}

\begin{enumerate}
\item \textbf{公式1}：关节空间动力学（无外部力）
\begin{equation}
\tau = M(\theta)\ddot{\theta} + h(\theta, \dot{\theta})
\end{equation}

\item \textbf{公式3}：关节空间动力学（含外部力）
\begin{equation}
\tau = M(\theta)\ddot{\theta} + h(\theta, \dot{\theta}) - J^T(\theta)F_{\text{ext}}
\end{equation}

\item \textbf{公式2}：任务空间动力学（从公式1或公式3转换而来）
\begin{equation}
F = \Lambda(\theta)\dot{V} + \eta(\theta, V)
\end{equation}
\end{enumerate}

\subsection{关键要点}

\begin{enumerate}
\item \textbf{公式1和公式3的区别}：
\begin{itemize}
\item 公式3多了一项$-J^T(\theta)F_{\text{ext}}$，表示外部力的影响
\item 公式1是公式3在$F_{\text{ext}} = 0$时的特殊情况
\end{itemize}

\item \textbf{公式2与公式1、公式3的关系}：
\begin{itemize}
\item 公式2是从公式1或公式3通过坐标变换得到的
\item 从关节空间转换到任务空间
\item 使用雅可比矩阵$J(\theta)$进行转换
\end{itemize}

\item \textbf{应用场景}：
\begin{itemize}
\item 公式1：自由空间运动
\item 公式3：接触任务
\item 公式2：任务空间控制
\end{itemize}
\end{enumerate}

\subsection{转换关系总结}

\textbf{从关节空间到任务空间}：
\begin{align}
\Lambda(\theta) &= J^{-T}(\theta)M(\theta)J^{-1}(\theta) \\
\eta(\theta, V) &= J^{-T}(\theta)h(\theta, J^{-1}(\theta)V) - \Lambda(\theta)\dot{J}(\theta)J^{-1}(\theta)V \\
V &= J(\theta)\dot{\theta} \\
\dot{V} &= \dot{J}(\theta)\dot{\theta} + J(\theta)\ddot{\theta}
\end{align}

\textbf{从任务空间到关节空间}：
\begin{align}
\tau &= J^T(\theta)F \\
\dot{\theta} &= J^{-1}(\theta)V \\
\ddot{\theta} &= J^{-1}(\theta)(\dot{V} - \dot{J}(\theta)J^{-1}(\theta)V)
\end{align}

\end{document}

